\section{Introduction}

Programming languages are not culturally neutral objects. Their grammars,
libraries, diagnostics, examples, documentation, and naming conventions all
encode assumptions about who is expected to read, write, and teach programs.
In contemporary programming practice, those assumptions are usually organized
around English. Even when an implementation offers translated documentation or
localized error messages, the language surface itself often remains built from
English words such as \code{if}, \code{while}, \code{return}, and \code{print}.

\plat{} explores a different design position. It is a small interpreted
language where the computational core is deliberately ordinary, but the visible
surface of the language is organized around Limburgian identity. Variables are
introduced with \code{loat}; functions with \code{funksie}; conditionals with
\code{es} and \code{angesj}; loops with \code{zolang} and \code{veur}; returns
with \code{trok}; and the tiny standard library includes \code{aafdrokke} for
printing and \code{invuier} for reading one line of input. The project also
contains English and Limburgian documentation paths and supports localized
diagnostics through command-line language tags such as \code{--lang li}.

The central claim of this paper is modest. \plat{} is not presented as a new
type system, runtime model, optimization strategy, or semantic foundation.
Instead, it is treated as an exploratory artifact: a concrete design case study
in what happens when regional-language identity is made a first-class
constraint in the design of a programming language. The conventionality of the
interpreter-language mechanisms is part of the method. By keeping the language
small and familiar, the project makes the cultural and linguistic layer more
visible as an object of study.

This draft makes four contributions:

\begin{enumerate}
  \item a concrete design case study of a Limburgian-themed programming
        language;
  \item a discussion of tradeoffs in regional-language keyword and diagnostic
        design;
  \item an analysis of authenticity versus usability, especially ASCII
        compatibility; and
  \item a reusable set of design questions for culturally localized programming
        languages.
\end{enumerate}

% TODO: Add a stable repository URL, release tag, archive DOI, and commit hash.
% TODO: Add a short methodology paragraph explaining how repository artifacts,
% documentation, examples, and implementation files were analyzed.
